\documentclass[11pt]{article}
\usepackage{amsmath,amsfonts,amssymb,amsthm}
\usepackage{graphicx}
\usepackage{algorithm}
\usepackage{algorithmic}
\usepackage{float}
\usepackage{hyperref}

\title{The Random Feature Method for Time-Dependent Problems}
\author{Jingrun Chen, Weinan E, and Yixin Luo}
\date{}

\begin{document}
\maketitle

\section{Introduction}

Time-dependent partial differential equations (PDEs) are widely used for modeling dynamic physical systems but pose challenges including complex geometries, mesh generation, and high dimensionality. Traditional numerical methods like finite difference, finite element, and spectral methods have been successful but face limitations. Meanwhile, deep learning approaches have gained attention for solving PDEs without mesh generation, but often lack reliable error control and clear convergence properties.

This paper introduces the Space-Time Random Feature Method (ST-RFM), which combines the advantages of both traditional and machine learning approaches. ST-RFM aims to create a mesh-free method with spectral accuracy in both space and time, particularly useful for complex geometries. The approach extends the Random Feature Method (RFM) to time-dependent problems through space-time partitioning and introduces two alternatives for constructing random feature functions: Space-Time Concatenation (STC) and Separation of Variables (SoV).

\section{Methodology}

\subsection{Background: Random Feature Method for Static Problems}

The random feature method (RFM) is first introduced for static boundary-value problems of the form:
\begin{equation}
\begin{cases}
Lu(x) = f(x), & x \in \Omega \\
Bu(x) = g(x), & x \in \partial\Omega
\end{cases}
\end{equation}
where $L$ and $B$ are differential and boundary operators, respectively.

The RFM methodology includes:
\begin{enumerate}
    \item Domain Decomposition: The domain $\Omega$ is decomposed into $N$ subdomains $\{\Omega_i\}$.
    
    \item Partition of Unity (PoU): For each subdomain $\Omega_i$, a PoU function $\psi_i$ with support on $\Omega_i$ is constructed. Two common PoU functions are:
    \begin{align}
        \psi_a(x) &= I_{[-1,1]}(x) \\
        \psi_b(x) &= I_{[-\frac{5}{4},-\frac{3}{4}]}(x)\frac{1 + \sin(2\pi x)}{2} + I_{[-\frac{3}{4},\frac{3}{4}]}(x) + I_{[\frac{3}{4},\frac{5}{4}]}(x)\frac{1 - \sin(2\pi x)}{2}
    \end{align}
    where $\psi_a$ is discontinuous while $\psi_b$ is continuously differentiable.
    
    \item Random Features: For each subdomain $\Omega_i$, $J_n$ random features $\varphi_{ij}$ are constructed:
    \begin{equation}
        \varphi_{ij}(x) = \sigma(k_{ij}^T l_i(x) + b_{ij})
    \end{equation}
    where $\sigma$ is an activation function (like tanh), $k_{ij}$ and $b_{ij}$ are randomly sampled parameters from $U(-R_m, R_m)$, and $l_i(x)$ is a linear transformation mapping inputs from $\Omega_i$ to $[-1,1]^{d_x}$.
    
    \item Approximate Solution: The solution is approximated as:
    \begin{equation}
        u_M(x) = \sum_{i=1}^{N} \psi_i(x) \sum_{j=1}^{J_n} u_{ij} \varphi_{ij}(x)
    \end{equation}
    
    \item Least-Squares Formulation: The coefficients $\{u_{ij}\}$ are determined by minimizing:
    \begin{equation}
        \text{Loss}(\{u_{i,j,k}\}) = \sum_{i=1}^{N} \sum_{q=1}^{Q} \| \alpha^p_{i,q} (Lu_M(x_{i,q}) - f(x_{i,q})) \|_2^2 + \sum_{x_{i,e} \in \partial\Omega} \| \alpha^b_{i,e} (Bu_M(x_{i,e}) - g(x_{i,e})) \|_2^2
    \end{equation}
    where $\alpha^p_{i,q}$ and $\alpha^b_{i,e}$ are rescaling parameters.
\end{enumerate}

\subsection{Space-Time Random Feature Method (ST-RFM)}

For time-dependent PDEs of the form:
\begin{equation}
\begin{cases}
Lu(x,t) = f(x,t), & x,t \in \Omega \times [0,T] \\
Bu(x,t) = g(x,t), & x,t \in \partial\Omega \times [0,T] \\
Iu(x,0) = h(x), & x \in \Omega
\end{cases}
\end{equation}
where $I$ is the initial operator, the ST-RFM extends the static RFM as follows:

\begin{enumerate}
    \item Space-Time Decomposition: The spatial domain $\Omega$ is decomposed into $N_x$ subdomains, and the temporal interval $[0,T]$ into $N_t$ subdomains.
    
    \item Space-Time PoU: A space-time PoU is constructed as:
    \begin{equation}
        \psi_{i_x,i_t}(x,t) = \psi_{i_x}(x) \psi_{i_t}(t)
    \end{equation}
    
    \item Space-Time Random Features: Two options are presented:
    \begin{itemize}
        \item Space-Time Concatenation (STC):
        \begin{equation}
            \varphi_{i_x,i_t,j}(x,t) = \sigma((k^x_{i_x,i_t,j})^T l_{i_x}(x) + k^t_{i_x,i_t,j} l_{i_t}(t) + b_{i_x,i_t,j})
        \end{equation}
        
        \item Separation of Variables (SoV):
        \begin{equation}
            \varphi_{i_x,i_t,j}(x,t) = \sigma((k^x_{i_x,i_t,j})^T l_{i_x}(x) + b^x_{i_x,i_t,j}) \sigma(k^t_{i_x,i_t,j} l_{i_t}(t) + b^t_{i_x,i_t,j})
        \end{equation}
    \end{itemize}
    
    \item Approximate Solution: The solution is approximated as:
    \begin{equation}
        u_M(x,t) = \sum_{i_x=1}^{N_x} \sum_{i_t=1}^{N_t} \psi_{i_x}(x) \psi_{i_t}(t) \sum_{j=1}^{J_n} u_{i_x,i_t,j} \varphi_{i_x,i_t,j}(x,t)
    \end{equation}
    
    \item Least-Squares Formulation: The coefficients are determined by minimizing a loss function that includes terms for the PDE, boundary conditions, and initial conditions.
\end{enumerate}

\subsection{Block Time-Marching Strategy}

The paper also discusses an alternative block time-marching strategy, which solves the time-dependent PDE sequentially in time blocks, using the solution at the end of one block as the initial condition for the next. This is detailed in the following algorithm:

\begin{algorithmic}[1]
\STATE $\tilde{A}_1 = [(\mathbf{\Phi}_{1,0})^T, (L_1)^T]^T$, $\tilde{b}_1 = [g^T, f_1^T]^T$
\STATE $u^B_1 = \min_u \|\tilde{A}_1 u - \tilde{b}_1\|_2$
\FOR{$i=2,3,\ldots,N_b$}
    \STATE $\tilde{A}_i = [(\mathbf{\Phi}_{i,0})^T, (L_i)^T]^T$, $\tilde{b}_i = [(u^B_{i-1})^T (\mathbf{\Phi}_{i-1,1})^T, f_i^T]^T$
    \STATE $u^B_i = \min_u \|\tilde{A}_i u - \tilde{b}_i\|_2$
\ENDFOR
\STATE $u^B = [(u^B_1)^T, (u^B_2)^T, \ldots, (u^B_{N_b})^T]^T$
\end{algorithmic}

\section{Theoretical Findings}

\subsection{Approximation Properties}

Both the STC and SoV formulations possess universal approximation properties:

\begin{theorem}[STC Approximation]
Let $\sigma$ be any continuous sigmoidal function. Then finite sums of the form:
\begin{equation}
G_C(x,t) = \sum_{j=1}^{N} u_j \sigma((k^x_j)^T x + k^t_j t + b_j)
\end{equation}
are dense in $C(\Omega \times [0,T])$.
\end{theorem}

\begin{corollary}
When using tanh as the activation function, the finite sums are dense in $C(\Omega \times [0,T])$.
\end{corollary}

\begin{theorem}[SoV Approximation]
For any continuous sigmoidal function $\sigma$, the finite sums:
\begin{equation}
G_S(x,t) = \sum_{j=1}^{N} u_j \sigma((k^x_j)^T x + b^x_j) \sigma(k^t_j t + b^t_j)
\end{equation}
are dense in $C(\Omega \times [0,T])$.
\end{theorem}

\subsection{Error Analysis}

The paper provides a detailed error analysis comparing the ST-RFM with the block time-marching strategy:

\begin{theorem}[Lower Bound for Block Time-Marching]
Under appropriate conditions, the error of the block time-marching strategy is bounded below by:
\begin{equation}
E_{\Delta u}[\|u^B_M(x,t) - u_e(x,t)\|_{L^2}] \geq \beta \sigma \sqrt{N_t} \phi(N_t) - \epsilon\sqrt{N_t}
\end{equation}
where $\phi(x) = 
\begin{cases}
1, & \text{if } \max(|\lambda_k|) \leq 1 \text{ and } 1 \not\in \{\lambda_k\} \\
x, & \text{if } \max(|\lambda_k|) = 1 \text{ and } 1 \in \{\lambda_k\} \\
|\lambda_m|^{x/2}, & \text{if } \max(|\lambda_k|) > 1
\end{cases}$
\end{theorem}

This shows that the error grows exponentially with the number of time blocks when $\max(|\lambda_k|) > 1$.

\begin{theorem}[Upper Bound for ST-RFM]
Under appropriate conditions, the error of the ST-RFM is bounded above by:
\begin{equation}
E_{\Delta u}[\|u^S(x,t) - u_e(x,t)\|_{L^2}] \leq \beta \Delta \sqrt{N_t} + \epsilon\sqrt{N_t}
\end{equation}
\end{theorem}

This indicates that the error in ST-RFM grows only sublinearly with the number of time subdomains, which is a significant improvement over the block time-marching strategy.

\section{Numerical Examples and Results}

\subsection{One-Dimensional Problems}

The paper presents numerical experiments for several one-dimensional problems:

\subsubsection{Heat Equation}
\begin{equation}
\begin{cases}
\partial_t u(x,t) - \beta^2 \partial^2_x u(x,t) = 0, & x \in [x_0, x_1], t \in [0,T] \\
u(x_0, t) = g_1(t), u(x_1, t) = g_2(t), & t \in [0,T] \\
u(x, 0) = h(x), & x \in [x_0, x_1]
\end{cases}
\end{equation}

Results show:
\begin{itemize}
    \item Both STC and SoV achieve spectral accuracy with L$^1$ error < 4.5e-6
    \item Convergence tests demonstrate spectral accuracy with respect to $N_t$, $J_n$, and the number of collocation points
    \item STC performs better than SoV for this problem
    \item ST-RFM maintains consistent accuracy across time subdomains, while block time-marching shows exponentially growing error
\end{itemize}

\subsubsection{Wave Equation}
\begin{equation}
\begin{cases}
\partial^2_t u(x,t) - \beta^2 \partial^2_x u(x,t) = 0, & x, t \in [x_0, x_1] \times [0,T] \\
u(x_0, t) = u(x_1, t) = 0, & t \in [0,T] \\
u(x, 0) = g_1(x), \partial_t u(x, 0) = g_2(x), & x \in [x_0, x_1]
\end{cases}
\end{equation}

Results show:
\begin{itemize}
    \item Both approaches achieve L$^1$ error < 2.5e-8
    \item SoV performs better than STC for this problem
    \item The ST-RFM maintains consistent accuracy, while block time-marching shows exponentially growing error
\end{itemize}

\subsubsection{Schrödinger Equation}
\begin{equation}
\begin{cases}
i\partial_t \psi(x,t) + 0.5 \Delta \psi(x,t) = 0, & x, t \in [x_0, x_1] \times [0,T] \\
\psi(x, 0) = g(x), & x \in [x_0, x_1] \\
\psi(x_0, t) = \psi(x_1, t), & t \in [0,T] \\
\partial_x \psi(x_0, t) = \partial_x \psi(x_1, t), & t \in [0,T]
\end{cases}
\end{equation}

Results show:
\begin{itemize}
    \item Both approaches achieve L$^1$ error < 2.0e-9
    \item SoV performs better than STC for this problem
    \item ST-RFM shows slower error growth compared to the exponential growth in block time-marching
\end{itemize}

\subsection{Two-Dimensional Problems}

The paper also presents results for two-dimensional membrane vibration problems:

\subsubsection{Membrane Vibration over Simple Geometry}
\begin{equation}
\begin{cases}
\partial_{tt} u(x,y,t) - \beta^2 \Delta u(x,y,t) = 0, & (x,y), t \in \Omega \times [0,T] \\
u(x,y,0) = \phi(x,y), \partial_t u(x,y,0) = \psi(x,y), & (x,y) \in \Omega \\
u(x,y,t) = 0, & (x,y) \in \partial\Omega \times [0,T]
\end{cases}
\end{equation}

Results demonstrate:
\begin{itemize}
    \item Spectral accuracy in both spatial and temporal directions
    \item Similar performance between STC and SoV
\end{itemize}

\subsubsection{Membrane Vibration over Complex Geometry}

The paper demonstrates the method's capability to handle complex geometries by solving a membrane vibration problem over an irregular domain. Without an exact solution available, numerical convergence is shown by:
\begin{itemize}
    \item Consistent convergence toward a reference solution as degrees of freedom increase
    \item Robustness across different parameter settings
\end{itemize}

\section{Implications and Advantages}

\subsection{Comparison between STC and SoV}

The numerical experiments reveal important insights about the two random feature constructions:
\begin{itemize}
    \item STC performs better for diffusion-dominated problems (heat equation)
    \item SoV performs better for wave-like equations (wave and Schrödinger equations)
    \item Both approaches are theoretically sound with universal approximation properties
\end{itemize}

\subsection{Advantages over Block Time-Marching}

The ST-RFM demonstrates significant advantages over the block time-marching strategy:
\begin{itemize}
    \item Error growth is sublinear with respect to the number of time subdomains, compared to exponential growth in block time-marching
    \item Theoretical error bounds are confirmed by numerical results
    \item More stable for long-time simulations
\end{itemize}

\subsection{Mesh-Free Advantages}

The ST-RFM preserves the mesh-free nature of machine learning approaches while maintaining the accuracy of traditional methods:
\begin{itemize}
    \item Only requires collocation points, not structured meshes
    \item Particularly valuable for problems with complex geometries
    \item Spectral accuracy in both space and time for problems with sufficient regularity
\end{itemize}

\subsection{Computational Efficiency}

The method offers computational advantages:
\begin{itemize}
    \item Formulates a convex optimization problem (linear least-squares)
    \item Avoids the expensive training procedures of deep neural networks
    \item Provides reliable error control and clear convergence properties
\end{itemize}

\end{document} 